\chapter{Running the programs}\label{ch:programs}
\sloppy

This chapter lists every program and script that comes with \ac{} and
says how to start each one: from which directory, with which command and
options, both with a build from the sources and with the container image.
Table~\ref{tab:programs} is the index. What the programs compute is in
Chapter~\ref{ch:workflow} and in the tutorials; their options in the
input file are in Chapter~\ref{ch:inputs}.

\begin{table}[ht]
\centering\small
\caption{The programs and scripts of \ac{}.}\label{tab:programs}
\begin{tabular}{lp{0.36\textwidth}r}
\toprule
program & what it does & section \\
\midrule
\multicolumn{3}{l}{\emph{The three phases}} \\
\prog{thermal\_anharmonicity\_fase1} & resonances, coefficients, band list & \ref{sec:run-fase1} \\
\prog{thermal\_anharmonicity\_fase2a} & polyads of the starting states & \ref{sec:run-fase2a} \\
\prog{thermal\_anharmonicity\_fase2b} & diagonalisation, transitions, histograms & \ref{sec:run-fase2b} \\
\midrule
\multicolumn{3}{l}{\emph{Tools (\file{tools/}, Python)}} \\
\prog{gaussian2ac.py} & input files from a Gaussian log & \ref{sec:run-gaussian2ac} \\
\prog{anh\_sample.py} & starting states for $T>0$, density of states & \ref{sec:run-sample} \\
\prog{anh\_jobs.py} & many starting states split into jobs & \ref{sec:run-jobs} \\
\prog{anh\_assemble.py} & spectra from the histograms of the jobs & \ref{sec:run-assemble} \\
\prog{anh\_hopball.py} & largest possible size of the polyads & \ref{sec:run-hopball} \\
\midrule
\multicolumn{3}{l}{\emph{Checks}} \\
\prog{tests/selftest.sh} & Tutorial~1 as a check of a build & \ref{sec:run-selftest} \\
\prog{anh-selftest} & the same, in the container image & \ref{sec:run-selftest} \\
\midrule
\multicolumn{3}{l}{\emph{Scripts of the examples (\file{examples/})}} \\
\prog{pyrene/run.sh} & Tutorial~1, the three phases & \ref{sec:run-examples} \\
\prog{pyrene/plot\_spectrum.py} & strongest lines and figure of Tutorial~1 & \ref{sec:run-examples} \\
\prog{pyrene-thermal/run\_sample.sh} & Tutorial~3, step~1 & \ref{sec:run-examples} \\
\prog{pyrene-thermal/thermal\_energy.py} & thermal energy distribution & \ref{sec:run-examples} \\
\bottomrule
\end{tabular}
\end{table}

The development repository also has \prog{tools/polyad\_overlap.py} and
\prog{tools/costmodel\_toy.py}, used to study the method; they are not
described here.

\section{Three ways to run them}\label{sec:run-ways}

In this chapter \file{CODEDIR} is the top directory of the code (the one
with \file{tools/}), and every command is given twice: first as it is typed with
a build from the sources, then with the container image, through a short
shell function \prog{anh} defined below. With the image, the tools are in
\nolinkurl{/opt/anharmonicaos/share/anharmonicaos/tools} (called \file{TOOLS}
below) and the examples in
\nolinkurl{/opt/anharmonicaos/share/anharmonicaos/examples} (also
\texttt{\$ANH\_EXAMPLES} inside the container).

\paragraph{Built from the sources.} The three programs are in
\file{CODEDIR} after \prog{make} (in \file{PREFIX/bin} after
\prog{make install}); put that directory in \texttt{PATH}, or give the
full path. The tools need the Python packages of Section~\ref{sec:python}:
activate their virtual environment first, then run
\texttt{python3 CODEDIR/tools/NAME.py}.
\begin{lstlisting}
export PATH=CODEDIR:$PATH
source ~/venvs/anharmonicaos/bin/activate
\end{lstlisting}

\paragraph{Docker.} A container sees only the directories mounted into it.
The function below mounts the current directory at the same path inside
the container and starts there, passes the number of OpenMP threads
(the image sets it to 1 otherwise), and runs the command given:
\begin{lstlisting}
IMG=ghcr.io/gmulas/anharmonicaos:2026.11.1
anh() { docker run --rm -v "$PWD":"$PWD" -w "$PWD" \
          -e OMP_NUM_THREADS="${OMP_NUM_THREADS:-1}" "$IMG" "$@"; }
\end{lstlisting}
Mounting at the same path keeps absolute paths valid, for example the
links that \prog{anh\_jobs.py} makes to the phase~1 files. Every other
directory that the command reads or writes must be mounted too, with one
more \texttt{-v DIR:DIR} per directory: for example the parent directory,
if \file{inp} names \file{../t1/freqsfile}, or the phase~1 directory of
\prog{anh\_jobs.py}. Mounting a common parent once (\texttt{-v
\$HOME/runs:\$HOME/runs}) covers them all. With a Docker daemon running
as root, add \texttt{-{}-user \$(id -u):\$(id -g)}, or the files written
belong to root; with rootless Docker and with Docker Desktop do not add
it (Section~\ref{sec:container}).

\paragraph{Apptainer and Singularity.} With the \file{.sif} file of
Section~\ref{sec:container}:
\begin{lstlisting}
SIF=$HOME/anharmonicaos-2026.11.1.sif
anh() { apptainer exec --env OMP_NUM_THREADS="${OMP_NUM_THREADS:-1}" "$SIF" "$@"; }
\end{lstlisting}
Use \prog{singularity} where \prog{apptainer} is not installed
(\texttt{-{}-env} needs Apptainer 1.0 or Singularity 3.6 or later). Apptainer
starts in the current directory, runs as you, and by default mounts the
home directory, the current directory and \file{/tmp} (the system
configuration can change this). Other file systems, for example a scratch
area, must be added with \texttt{-{}-bind}, as in
\texttt{apptainer exec -{}-bind /scratch,/archive \ldots}. Inside a Slurm
job, add \texttt{-{}-cleanenv} when the command is \prog{mpirun} inside the
container: otherwise it reads the task count of the job and starts a
single process. Running one container per MPI process, across nodes, is
described in Section~\ref{sec:container}.

\paragraph{Messages.} All the programs write their log on the standard
output and errors and warnings on the standard error: redirect the first
to a file (\texttt{> fase1.log}) and keep the second on the terminal, or
redirect both. Each program writes to \file{out.oct} every option it read,
with the value used.

\section{\prog{thermal\_anharmonicity\_fase1}}\label{sec:run-fase1}

\paragraph{Where.} In the directory that contains \file{inp} and the
input files it names (\file{freqsfile}, \file{phi3file}, \file{phi4file},
\file{dmu2file}, Chapter~\ref{ch:inputs}).

\paragraph{Command.} It takes no arguments; the threshold tests use
OpenMP threads.
\begin{lstlisting}
OMP_NUM_THREADS=4 thermal_anharmonicity_fase1 > fase1.log
OMP_NUM_THREADS=4 anh thermal_anharmonicity_fase1 > fase1.log
\end{lstlisting}

\paragraph{Result.} \file{resonanceslist}, \file{dunhamparams},
\file{bandlist}, \file{resonancecandidates} and \file{dumpedparams}, in the
same directory (Chapter~\ref{ch:outputs}). Phase~1 alone can also be built
without MPI, PETSc or PARI (Section~\ref{sec:checks}).

\section{\prog{thermal\_anharmonicity\_fase2a}}\label{sec:run-fase2a}

\paragraph{Where.} In the directory of phase~1, after it, with the file of
starting states named by \opt{modepopfile} (\file{quantafile} in the
tutorials).

\paragraph{Command.} One process, no arguments:
\begin{lstlisting}
thermal_anharmonicity_fase2a > fase2a.log
anh thermal_anharmonicity_fase2a > fase2a.log
\end{lstlisting}

\paragraph{Result.} One directory \file{dir\_n00000}, \file{dir\_n00001},
\ldots{} per starting state, with the Hamiltonian matrices of its polyads,
and \file{polyadsizes}, \file{processedquantas}, \file{processeddirs}.
With \opt{sizesonly}~$=1$, only \file{growthsizes}.

\section{\prog{thermal\_anharmonicity\_fase2b}}\label{sec:run-fase2b}

\paragraph{Where.} In the same directory, after phase~2a.

\paragraph{Command.} It runs with MPI and takes no arguments. With the
distributed eigensolvers (\opt{eigensolver}~$=1$, 2, 3), give the number
of processes, at most the physical cores:
\begin{lstlisting}
mpirun -np 4 thermal_anharmonicity_fase2b > fase2b.log
anh mpirun -np 4 thermal_anharmonicity_fase2b > fase2b.log
\end{lstlisting}
With the spectral method (\opt{spectral}~$=1$, \opt{eigensolver}~$=0$) use
one process and OpenMP threads; \texttt{-{}-bind-to none} lets the threads
use all the cores:
\begin{lstlisting}
OMP_NUM_THREADS=8 mpirun -np 1 --bind-to none thermal_anharmonicity_fase2b > fase2b.log
OMP_NUM_THREADS=8 anh mpirun -np 1 --bind-to none thermal_anharmonicity_fase2b > fase2b.log
\end{lstlisting}
In a Slurm job, \prog{srun} can replace \prog{mpirun} for a build from
the sources; with Apptainer, add \texttt{-{}-cleanenv} to \prog{anh}
(Section~\ref{sec:run-ways}).

\paragraph{Result.} The histograms \file{hist\_*} and one log per starting
state, \file{proclog\_n*}.

\section{\prog{gaussian2ac.py}}\label{sec:run-gaussian2ac}

\paragraph{Where.} In the directory where the input files of \ac{} should
be written, or anywhere with \texttt{-{}-outdir}. The Gaussian log must come
from a \texttt{freq=(anharm,PrintDerivatives)} job
(Chapter~\ref{ch:tut-gaussian}).

\paragraph{Command.}
\begin{lstlisting}
python3 CODEDIR/tools/gaussian2ac.py molecule.log [--outdir DIR]
anh python3 TOOLS/gaussian2ac.py molecule.log [--outdir DIR]
\end{lstlisting}
\texttt{-{}-phinoisethres} and \texttt{-{}-dmunoisethres} change the
thresholds below which force constants and dipole derivatives are taken
as numerical noise; \texttt{-{}-no-dmu3} skips \file{dmu3file}.

\paragraph{Result.} \file{freqsfile}, \file{phi3file}, \file{phi4file},
\file{dmu2file} and, when the log has them, \file{dmu3file}.

\section{\prog{anh\_sample.py}}\label{sec:run-sample}

\paragraph{Where.} Anywhere: it reads \file{inp} and \file{freqsfile}
from the paths given, and writes in the current directory.

\paragraph{Command.} \texttt{-{}-enum} and \texttt{-{}-emax} are required
(Section~\ref{sec:thermal-sample}):
\begin{lstlisting}
python3 CODEDIR/tools/anh_sample.py --inp inp --enum 1000 --emax 2500 \
  --per-bin 2 --blocks 1 --prefix s
anh python3 TOOLS/anh_sample.py --inp inp --enum 1000 --emax 2500 \
  --per-bin 2 --blocks 1 --prefix s
\end{lstlisting}
\texttt{-{}-walkers N --procs N} runs N Wang--Landau walkers in parallel.

\paragraph{Result.} \file{PREFIX\_states.dat}, \file{PREFIX\_dos.dat} and,
with \texttt{-{}-per-bin}, the blocks \file{PREFIX\_quanta\_block000},
\ldots{} in the \file{quantafile} format.

\section{\prog{anh\_jobs.py}}\label{sec:run-jobs}

\paragraph{Where.} Anywhere: \texttt{-{}-base} is the directory where
phase~1 was run, \texttt{-{}-out} the directory of the jobs, which must not
exist (Section~\ref{sec:thermal-jobs}).

\paragraph{Command.}
\begin{lstlisting}
python3 CODEDIR/tools/anh_jobs.py --base ../t1 --profile laptop --out jobs \
  --states s_quanta_block000 --Ehiststep 100 --Emax 2500
anh python3 TOOLS/anh_jobs.py --base ../t1 --profile laptop --out jobs \
  --states s_quanta_block000 --Ehiststep 100 --Emax 2500
\end{lstlisting}
With Docker, mount the phase~1 directory too (here the parent directory:
\texttt{-v "\$(dirname \$PWD)":"\$(dirname \$PWD)"}), since the jobs link
to its files.

\paragraph{Running the jobs.} With a build from the sources:
\begin{lstlisting}
ANHBIN=CODEDIR sh jobs/submit_all.sh
\end{lstlisting}
With the container and a shell profile (\texttt{laptop}), run the whole
script inside it: the programs are in its \texttt{PATH}.
\begin{lstlisting}
anh sh jobs/submit_all.sh
\end{lstlisting}
With Apptainer and a Slurm profile, \file{submit\_all.sh} runs on the
host and each job starts the container. Write the container in the
\texttt{mpirun} entry of the profile, with the threads of the profile,
and leave \texttt{bindir} empty:
\begin{lstlisting}
"mpirun": "apptainer exec --cleanenv --env OMP_NUM_THREADS=8 /path/anharmonicaos-2026.11.1.sif mpirun -np 1 --bind-to none",
"bindir": ""
\end{lstlisting}
Add \texttt{-{}-bind} for the file systems of the jobs and of the phase~1
directory if they are not mounted by default.

\section{\prog{anh\_assemble.py}}\label{sec:run-assemble}

\paragraph{Where.} Anywhere: \texttt{-{}-jobs} is the directory written by
\prog{anh\_jobs.py} (or \texttt{-{}-dirs} lists phase~2b directories), and
\texttt{-{}-out} the directory of the results
(Section~\ref{sec:thermal-assemble}).

\paragraph{Command.}
\begin{lstlisting}
python3 CODEDIR/tools/anh_assemble.py --jobs jobs --out res --T 0 100 200 --lng s_dos.dat
anh python3 TOOLS/anh_assemble.py --jobs jobs --out res --T 0 100 200 --lng s_dos.dat
\end{lstlisting}
\texttt{-{}-plot} also writes figures (matplotlib).

\paragraph{Result.} In \file{res/}: the summed \file{hist\_*},
\file{spectra\_by\_energy.dat}, \file{spectra\_by\_anh\_energy.dat}, and,
with \texttt{-{}-T}, \file{thermal\_T<T>.dat}.

\section{\prog{anh\_hopball.py}}\label{sec:run-hopball}

\paragraph{Where.} In the directory of phase~1 (it reads
\file{resonanceslist}, and \file{bandlist} with \texttt{-{}-bandlist}).

\paragraph{Command.}
\begin{lstlisting}
python3 CODEDIR/tools/anh_hopball.py resonanceslist --hops 4 \
  --quanta quantafile --bandlist bandlist
anh python3 TOOLS/anh_hopball.py resonanceslist --hops 4 \
  --quanta quantafile --bandlist bandlist
\end{lstlisting}

\paragraph{Result.} On the terminal, the number of states within 0 to
$H$ steps (Chapter~\ref{ch:workflow}).

\section{The checks}\label{sec:run-selftest}

\prog{tests/selftest.sh} checks a build from the sources, and
\prog{anh-selftest} the container image. Both run Tutorial~1 in a
temporary directory, in about a minute, with N MPI processes (default 2,
at most the physical cores), and end with \texttt{OK}:
\begin{lstlisting}
CODEDIR/tests/selftest.sh CODEDIR 4         # from anywhere
docker run --rm ghcr.io/gmulas/anharmonicaos:2026.11.1 anh-selftest 4
apptainer exec anharmonicaos-2026.11.1.sif anh-selftest 4
\end{lstlisting}
The first argument of \prog{selftest.sh} is the directory of the
programs. The numerical checks of \file{tests/checks} are described in
Section~\ref{sec:checks} and in \file{tests/checks/README.md}.

\section{The scripts of the examples}\label{sec:run-examples}

\paragraph{\prog{examples/pyrene/run.sh BINDIR [NPROC]}.} Runs Tutorial~1
(the three phases, phase~2b with NPROC processes, default 4) in the
directory of the example, and keeps a copy of \file{out.oct} after each
phase. Copy the directory first, so the original stays clean:
\begin{lstlisting}
cp -r CODEDIR/doc/manual/examples/pyrene t1 && cd t1 && sh run.sh CODEDIR 4
anh cp -r /opt/anharmonicaos/share/anharmonicaos/examples/pyrene t1
cd t1 && anh sh run.sh /opt/anharmonicaos/bin 4
\end{lstlisting}

\paragraph{\prog{examples/pyrene/plot\_spectrum.py [RUNDIR] [NLINES]}.}
Prints the NLINES strongest lines of a finished run of Tutorial~1 and,
with matplotlib, draws \file{pyrene-spectrum.png} in RUNDIR. \file{run.sh} calls it
at the end.
\begin{lstlisting}
python3 plot_spectrum.py . 15
anh python3 plot_spectrum.py . 15
\end{lstlisting}

\paragraph{\prog{examples/pyrene-thermal/run\_sample.sh CODEDIR}.} Step~1
of Tutorial~3: runs \prog{anh\_sample.py} on the inputs of Tutorial~1, in
the current directory, in a few seconds.
\begin{lstlisting}
sh CODEDIR/doc/manual/examples/pyrene-thermal/run_sample.sh CODEDIR
anh sh /opt/anharmonicaos/share/anharmonicaos/examples/pyrene-thermal/run_sample.sh \
  /opt/anharmonicaos/share/anharmonicaos
\end{lstlisting}

\paragraph{\prog{examples/pyrene-thermal/thermal\_energy.py [FREQSFILE]
[--plot]}.} For a few temperatures, the mean vibrational energy and the
energies below which 99\,\% and 99.9\,\% of the population lie
(Section~\ref{sec:thermal-howhigh}); \texttt{-{}-plot} draws the figure.
\begin{lstlisting}
python3 thermal_energy.py freqsfile --plot
anh python3 thermal_energy.py freqsfile --plot
\end{lstlisting}
\fussy
